\documentclass[10pt,a4paper]{amsart}
\usepackage{amsmath,amssymb,amsthm,graphicx}
\usepackage{verbatim}
\newenvironment{reference}{\comment}{\endcomment}
%\newenvironment{reference}{}{}
\newenvironment{slogan}{\comment}{\endcomment}
\newenvironment{history}{\comment}{\endcomment}

% For commutative diagrams we use Xy-pic
\usepackage[all]{xy}

% We use 2cell for 2-commutative diagrams.
\xyoption{2cell}
% Theorem environments.
%
\theoremstyle{plain}
\newtheorem{theorem}[subsection]{Theorem}
\newtheorem{proposition}[subsection]{Proposition}
\newtheorem{lemma}[subsection]{Lemma}

\theoremstyle{definition}
\newtheorem{definition}[subsection]{Definition}
\newtheorem{example}[subsection]{Example}
\newtheorem{exercise}[subsection]{Exercise}
\newtheorem{situation}[subsection]{Situation}

\theoremstyle{remark}
\newtheorem{remark}[subsection]{Remark}
\newtheorem{remarks}[subsection]{Remarks}

\numberwithin{equation}{subsection}

% Macros
%
\def\lim{\mathop{\mathrm{lim}}\nolimits}
\def\colim{\mathop{\mathrm{colim}}\nolimits}
\def\Spec{\mathop{\mathrm{Spec}}}
\def\Hom{\mathop{\mathrm{Hom}}\nolimits}
\def\Ext{\mathop{\mathrm{Ext}}\nolimits}
\def\SheafHom{\mathop{\mathcal{H}\!\mathit{om}}\nolimits}
\def\SheafExt{\mathop{\mathcal{E}\!\mathit{xt}}\nolimits}
\def\Sch{\mathit{Sch}}
\def\Mor{\mathop{\mathrm{Mor}}\nolimits}
\def\Ob{\mathop{\mathrm{Ob}}\nolimits}
\def\Sh{\mathop{\mathit{Sh}}\nolimits}
\def\NL{\mathop{N\!L}\nolimits}
\def\CH{\mathop{\mathrm{CH}}\nolimits}
\def\proetale{{pro\text{-}\acute{e}tale}}
\def\etale{{\acute{e}tale}}
\def\QCoh{\mathit{QCoh}}
\def\Ker{\mathop{\mathrm{Ker}}}
\def\Im{\mathop{\mathrm{Im}}}
\def\Coker{\mathop{\mathrm{Coker}}}
\def\Coim{\mathop{\mathrm{Coim}}}

% Boxtimes
%
\DeclareMathSymbol{\boxtimes}{\mathbin}{AMSa}{"02}

%
% Macros for moduli stacks/spaces
%
\def\QCohstack{\mathcal{QC}\!\mathit{oh}}
\def\Cohstack{\mathcal{C}\!\mathit{oh}}
\def\Spacesstack{\mathcal{S}\!\mathit{paces}}
\def\Quotfunctor{\mathrm{Quot}}
\def\Hilbfunctor{\mathrm{Hilb}}
\def\Curvesstack{\mathcal{C}\!\mathit{urves}}
\def\Polarizedstack{\mathcal{P}\!\mathit{olarized}}
\def\Complexesstack{\mathcal{C}\!\mathit{omplexes}}
% \Pic is the operator that assigns to X its picard group, usage \Pic(X)
% \Picardstack_{X/B} denotes the Picard stack of X over B
% \Picardfunctor_{X/B} denotes the Picard functor of X over B
\def\Pic{\mathop{\mathrm{Pic}}\nolimits}
\def\Picardstack{\mathcal{P}\!\mathit{ic}}
\def\Picardfunctor{\mathrm{Pic}}
\def\Deformationcategory{\mathcal{D}\!\mathit{ef}}
\usepackage{cleveref}
\crefname{originaldefinition}{original definition}{original definitions}
\Crefname{originaldefinition}{Original definition}{Original definitions}
\makeatletter
\newlabel{definition-locally-finite-presentation}{{03XP}{source}{}{sourcepointer-1}{}}
\newlabel{definition-locally-finite-presentation@cref}{{[originaldefinition][0][]03XP}{[1][0][]source}{}{}{}}
\makeatother
\crefname{originallemma}{original lemma}{original lemmas}
\Crefname{originallemma}{Original lemma}{Original lemmas}
\makeatletter
\newlabel{lemma-finite-presentation-finite-type}{{0464}{source}{}{sourcepointer-2}{}}
\newlabel{lemma-finite-presentation-finite-type@cref}{{[originallemma][0][]0464}{[1][0][]source}{}{}{}}
\makeatother
\crefname{originallemma}{original lemma}{original lemmas}
\Crefname{originallemma}{Original lemma}{Original lemmas}
\makeatletter
\newlabel{lemma-locally-finite-type-locally-noetherian}{{04ZK}{source}{}{sourcepointer-3}{}}
\newlabel{lemma-locally-finite-type-locally-noetherian@cref}{{[originallemma][0][]04ZK}{[1][0][]source}{}{}{}}
\makeatother
\crefname{originallemma}{original lemma}{original lemmas}
\Crefname{originallemma}{Original lemma}{Original lemmas}
\makeatletter
\newlabel{lemma-separated-over-separated}{{04ZH}{source}{}{sourcepointer-4}{}}
\newlabel{lemma-separated-over-separated@cref}{{[originallemma][0][]04ZH}{[1][0][]source}{}{}{}}
\makeatother
\crefname{originaldefinition}{original definition}{original definitions}
\Crefname{originaldefinition}{Original definition}{Original definitions}
\makeatletter
\newlabel{spaces-properties-definition-noetherian}{{03EA}{source}{}{sourcepointer-5}{}}
\newlabel{spaces-properties-definition-noetherian@cref}{{[originaldefinition][0][]03EA}{[1][0][]source}{}{}{}}
\makeatother
\title{A finite presentation morphism of algebraic spaces over a Noetherian base has Noetherian source}
\author{The Stacks Project Authors}
\date{stacks-project:spaces-morphisms.tex:local-snapshot-20261012; GNU FDL 1.2 or later}
\begin{document}
\maketitle
\setcounter{subsection}{1}
\noindent\textit{Editorial note.} All delivered mathematical text is the author's own, copied from the pinned Stacks Project chapter file spaces-morphisms.tex; no statement or proof step was rewritten, reordered or removed.\par\medskip
\noindent\textit{Editorial note.} The proof is a verification of the three conditions of the definition; each step names the lemma or definition it uses, and no step is left to the reader.\par\medskip
\noindent\textit{Editorial note.} The five references of the counted spans are recorded as non-load-bearing navigation pointers; the argument that consumes them stays inside the counted proof span.\par\medskip
\noindent\textit{Editorial note.} Preamble: the three native excerpt spans of sources/stacks-spaces-morphisms/preamble.tex (lines 8-13, 14-18 and 154-217). The chapter file itself carries no preamble text.\par\medskip
\noindent\textit{Allowed standard prerequisites (editorial).}\begin{itemize}
\item A locally finite type morphism to a locally Noetherian base has locally Noetherian source. Lemmas lemma-finite-presentation-finite-type and lemma-locally-finite-type-locally-noetherian, used at proof lines 5771-5774. Gives the first of the three defining conditions.
\item Finite presentation makes a morphism quasi-compact and quasi-separated. Definition definition-locally-finite-presentation, used at proof lines 5776-5777. Combined with quasi-compactness of Y gives quasi-compactness of X.
\item Noetherianness of an algebraic space in terms of the three conditions. Properties of Spaces, Definition spaces-properties-definition-noetherian, used at proof lines 5778-5786. Is the definition whose three conditions the proof verifies.
\item Separatedness over a quasi-separated base is stable. Lemma lemma-separated-over-separated, used at proof lines 5781-5782. Gives quasi-separatedness of X.
\end{itemize}\medskip
\ifdefined\hypertarget\hypertarget{sourcepointer-1}{}\fi
\noindent\textit{Source reference.} Original definition 03XP: focused-seeds/AGP/sources/stacks-snapshot-20261011/tags.txt, source lines 5156--5156. Navigation pointer into the Stacks Project; the standard fact it supplies is applied inside the counted proof span and the argument resting on it stays there.\par\medskip
\ifdefined\hypertarget\hypertarget{sourcepointer-2}{}\fi
\noindent\textit{Source reference.} Original lemma 0464: focused-seeds/AGP/sources/stacks-snapshot-20261011/tags.txt, source lines 5459--5459. Navigation pointer into the Stacks Project; the standard fact it supplies is applied inside the counted proof span and the argument resting on it stays there.\par\medskip
\ifdefined\hypertarget\hypertarget{sourcepointer-3}{}\fi
\noindent\textit{Source reference.} Original lemma 04ZK: focused-seeds/AGP/sources/stacks-snapshot-20261011/tags.txt, source lines 6468--6468. Navigation pointer into the Stacks Project; the standard fact it supplies is applied inside the counted proof span and the argument resting on it stays there.\par\medskip
\ifdefined\hypertarget\hypertarget{sourcepointer-4}{}\fi
\noindent\textit{Source reference.} Original lemma 04ZH: focused-seeds/AGP/sources/stacks-snapshot-20261011/tags.txt, source lines 6465--6465. Navigation pointer into the Stacks Project; the standard fact it supplies is applied inside the counted proof span and the argument resting on it stays there.\par\medskip
\ifdefined\hypertarget\hypertarget{sourcepointer-5}{}\fi
\noindent\textit{Source reference.} Original definition 03EA: focused-seeds/AGP/sources/stacks-snapshot-20261011/tags.txt, source lines 4449--4449. Navigation pointer into the Stacks Project; the standard fact it supplies is applied inside the counted proof span and the argument resting on it stays there.\par\medskip
\begin{lemma}
\label{lemma-finite-presentation-noetherian}
Let $S$ be a scheme.
Let $f : X \to Y$ be a morphism of algebraic spaces over $S$.
If $f$ is of finite presentation and $Y$ is Noetherian,
then $X$ is Noetherian.
\end{lemma}

\begin{proof}
Assume $f$ is of finite presentation and $Y$ Noetherian. By
Lemmas \ref{lemma-finite-presentation-finite-type} and
\ref{lemma-locally-finite-type-locally-noetherian}
we see that $X$ is locally Noetherian. As $f$ is quasi-compact
and $Y$ is quasi-compact we see that $X$ is quasi-compact.
As $f$ is of finite presentation it is quasi-separated (see
Definition \ref{definition-locally-finite-presentation})
and as $Y$ is Noetherian it is quasi-separated (see
Properties of Spaces,
Definition \ref{spaces-properties-definition-noetherian}).
Hence $X$ is quasi-separated by
Lemma \ref{lemma-separated-over-separated}.
Hence we have checked all three conditions of
Properties of Spaces,
Definition \ref{spaces-properties-definition-noetherian}
and we win.
\end{proof}

\end{document}
